\documentclass[10pt,a4paper]{amsart}
\usepackage[margin=19mm]{geometry}
\usepackage{amsmath,amssymb,mathtools}
\usepackage[scr=rsfs,cal=euler]{mathalfa}
\usepackage{xfrac}
\usepackage[unicode,hidelinks,bookmarks=false]{hyperref}
\newcommand{\ie}{i.e.\ }
\newcommand{\cf}{cf.\ }
\newcommand{\resp}{respectively\ }
\newcommand{\Subsection}{\subsection}
\providecommand{\nobreakdash}{\nobreak\hbox{-}}
\setlength{\emergencystretch}{2em}
\multlinegap0pt
\usepackage{amssymb}
\usepackage{mathtools}
\usepackage[shortlabels]{enumitem}
\newtheorem{lemma}{Lemma}
\newcommand{\NN}{\mathbb N}
\newcommand{\ZZ}{\mathbb Z}
\title{A Counterexample to the\ Huneke--Wiegand Conjecture}
\author{Son Pham}
\date{Source: arXiv:2609.07615v1 (CC BY 4.0)}
\begin{document}
\maketitle

\noindent\emph{Source and licence.} A Counterexample to the\ Huneke--Wiegand Conjecture. Version arXiv:2609.07615v1 (CC BY 4.0), distributed by arXiv under the Creative Commons Attribution 4.0 International licence, http://creativecommons.org/licenses/by/4.0/, as recorded in the arXiv licence metadata for this exact version and shown on the abstract page https://arxiv.org/abs/2609.07615v1. Full licence notice: \texttt{licenses/arxiv-2609.07615-CC-BY-4.0.txt}.\par\smallskip

\noindent\textit{Editorial scope.} Counted unit: the article's lemma lem:certificate (source0.tex lines 203--210). The article's complete original proof of it is delivered with it (source0.tex lines 211--241, one \texttt{proof} environment of 31 lines). No mathematics is written, repaired or re-derived: the statement and the proof are the article's own lines, in the article's own notation, and no retained line is substituted or re-flowed.  Retained uncounted as the source-local context the proof needs: lines 100--117; lines 584--584.  Nothing was omitted from any retained span, so the package carries no omission record.  No retained line contains a citation, so the wrapper carries no bibliography and no bibliography entry.  No retained line contains a figure environment, an image-inclusion command or drawing code, so no image is delivered and the package declares no asset dependency.  Editorial formatting only, in addition: an AMS \texttt{amsart} preamble that restates the article's own declarations and macros used by the retained lines, namely the environments \texttt{lemma} on separate counters, together with the standard packages those lines need.  The retained environments print in this document as Lemma 1 in order of appearance, whereas the article prints the same items with the numbering of its own full document.  The complete native source of the article as distributed in the arXiv e-print for this exact version is bundled unchanged as \texttt{sources/209644/source0.tex}.
\begin{lemma}\label{lem:membership}
The elements of $\Gamma$ through $181$ are
\begin{equation}\label{eq:membership}
\begin{split}
\Gamma\cap[0,181]=\{0\}&\cup[56,58]\cup[63,64]\cup[70,83]\cup\{87\}\\
&\cup[89,90]\cup\{93\}\cup[95,97]\cup[112,116]\\
&\cup[119,122]\cup[126,180].
\end{split}
\end{equation}
Every integer at least $182$ belongs to $\Gamma$. Thus the Frobenius number is $181$, the conductor is $182$, and
\begin{equation}\label{eq:gaps}
\begin{split}
\NN\setminus\Gamma={}&[1,55]\cup[59,62]\cup[65,69]\cup[84,86]\cup\{88\}\\
&\cup[91,92]\cup\{94\}\cup[98,111]\cup[117,118]\\
&\cup[123,125]\cup\{181\}.
\end{split}
\end{equation}
\end{lemma}

\section{Complete pair-sum certificate}\label{app:pairs}

\begin{lemma}[Finite colon certificate]\label{lem:certificate}
The following equalities of subsets of $\ZZ$ hold:
\begin{align}
 E&=B+\Gamma,\label{eq:E}\\
 B+B&=C,\label{eq:BB}\\
 E\cap(14+E)&=X+\Gamma.\label{eq:intersection}
\end{align}
\end{lemma}

\begin{proof}
Lemma~\ref{lem:membership} and the conductor give
\begin{equation}\label{eq:Einterval}
\begin{split}
E=B&\cup[112,116]\cup[119,122]\cup[126,166]\\
 &\cup[168,180]\cup[182,\infty).
\end{split}
\end{equation}
Direct addition gives
\[
 B+G=[112,116]\cup[119,122]\cup[126,166]\cup[168,180].
\]
The tail lies in $B+\Gamma$ because
\begin{align*}
 [182,222]&=56+[126,166],&223&=57+166,\\
 [224,236]&=56+[168,180],&237&=57+180,
\end{align*}
and addition by $56$ propagates the block $[182,237]$. This proves~\eqref{eq:E}; the reverse inclusion also follows at once because $E$ is closed under addition by $\Gamma$.

The equality~\eqref{eq:BB} is a finite addition table. Appendix~\ref{app:pairs} lists a representation of every element of $C$; enumerating all unordered pairs from $B$ produces no value outside $C$.

Finally, intersecting~\eqref{eq:Einterval} with its translate gives
\[
 E\cap(14+E)=X\cup[182,194]\cup[196,\infty).
\]
A further finite addition yields
\[
 X+G=[182,194]\cup[196,277].
\]
The last interval has length greater than $56$, so addition by $56$ supplies its entire tail. This proves~\eqref{eq:intersection}.
\end{proof}

\end{document}
